% GENERATED FILE. Edit canonical lessons, then run npm run generate:books.
% canonical-source-hash: fnv1a64:6f8d7117e7c3f43a
% canonical-lessons: ML-C398-simhasanam, ML-C398-kiritam, ML-C398-pattabhisekam, ML-C398-bharanadhikari, ML-C398-valuka

\chapter{The Throne and the Crown}
\label{ch:ml-the-throne-and-the-crown}
\hlchaptermodality{398}

\begin{chapteropening}
\textbf{By the end of this chapter:} \emph{I can say a throne, a crown, a coronation, a sovereign and to reign.}

Complete the last lesson of chapter 398: I can say a throne, a crown, a coronation, a sovereign and to reign.
\end{chapteropening}

\section[\texorpdfstring{\ml{സിംഹാസനം} (siṁhāsanaṁ)}{സിംഹാസനം (siṁhāsanaṁ)}]{\ml{സിംഹാസനം} (siṁhāsanaṁ) — a throne}

\label{lesson:ML-C398-simhasanam}



\begin{quote}
Before the new one: say the Malayalam for a princess, then the Malayalam for a palace.
\end{quote}

\subsection*{You'll want to know: \ml{സിംഹാസനം}}
\textbf{\ml{സിംഹാസനം}} — \emph{siṁhāsanaṁ} — \textquotedblleft{}a throne\textquotedblright{}.

\begin{quote}
\ml{രാജാവ്} \ml{സിംഹാസനത്തിൽ} \ml{ഇരുന്നു}.
\end{quote}

\begin{quote}
\emph{rājāvŭ siṁhāsanattil irunnu.}
\end{quote}

\begin{quote}
The king sat on the throne.
\end{quote}

Literally a lion seat: \ml{സിംഹം}, a lion, + \ml{ആസനം}, a seat.

\begin{cousinweb}
\textbf{\ml{സിംഹാസനം}} (\emph{siṁhāsanaṁ}) is borrowed from Sanskrit \emph{siṁhāsana}, \textquotedblleft{}a lion's seat, a king's seat, a throne\textquotedblright{} (Monier-Williams), from \emph{siṁha}, \textquotedblleft{}a lion\textquotedblright{} (CDIAL 13384), and \emph{āsana}, \textquotedblleft{}a seat\textquotedblright{} (CDIAL 1484).
\end{cousinweb}

\begin{grammarlens}[title={What you've built}]
One more word for telling what happened.
\end{grammarlens}

\subsection*{Your turn}
\begin{itemize}
\raggedright
  \item \emph{Say it:} \emph{siṁhāsanaṁ}
  \item \emph{Say it:} \emph{siṁhāsanaṁ}, once more
  \item \emph{Say it:} \emph{rājāvŭ siṁhāsanattil irunnu.} — The king sat on the throne.
  \item \emph{Recall:} say the Malayalam for a princess, then the Malayalam for a palace, then say \emph{siṁhāsanaṁ} again
\end{itemize}

\begin{tcolorbox}[breakable,colback=teal!4,colframe=teal!35!black,title={Before you move on}]
What does \ml{സിംഹാസനം} mean? (\textquotedblleft{}A throne\textquotedblright{}.) Say it once more.
\end{tcolorbox}
\section[\texorpdfstring{\ml{കിരീടം} (kirīṭaṁ)}{കിരീടം (kirīṭaṁ)}]{\ml{കിരീടം} (kirīṭaṁ) — a crown}

\label{lesson:ML-C398-kiritam}



\begin{quote}
Before the new one: say the Malayalam for a palace, then the Malayalam for a throne.
\end{quote}

\subsection*{You'll want to know: \ml{കിരീടം}}
\textbf{\ml{കിരീടം}} — \emph{kirīṭaṁ} — \textquotedblleft{}a crown\textquotedblright{}.

\begin{quote}
\ml{രാജാവിന്റെ} \ml{കിരീടം} \ml{സ്വർണ്ണം} \ml{കൊണ്ടുള്ളതായിരുന്നു}.
\end{quote}

\begin{quote}
\emph{rājāvinṟe kirīṭaṁ svarṇṇaṁ koṇṭuḷḷatāyirunnu.}
\end{quote}

\begin{quote}
The king's crown was made of gold.
\end{quote}

Also a title in sport: the winner takes the \ml{കിരീടം}.

\begin{cousinweb}
\textbf{\ml{കിരീടം}} (\emph{kirīṭaṁ}) is borrowed from Sanskrit \emph{kirīṭa}, \textquotedblleft{}a diadem, a crest\textquotedblright{} (CDIAL 3176).
\end{cousinweb}

\begin{grammarlens}[title={What you've built}]
One more word for telling what happened.
\end{grammarlens}

\subsection*{Your turn}
\begin{itemize}
\raggedright
  \item \emph{Say it:} \emph{kirīṭaṁ}
  \item \emph{Say it:} \emph{kirīṭaṁ}, once more
  \item \emph{Say it:} \emph{rājāvinṟe kirīṭaṁ svarṇṇaṁ koṇṭuḷḷatāyirunnu.} — The king's crown was made of gold.
  \item \emph{Recall:} say the Malayalam for a palace, then the Malayalam for a throne, then say \emph{kirīṭaṁ} again
\end{itemize}

\begin{tcolorbox}[breakable,colback=teal!4,colframe=teal!35!black,title={Before you move on}]
What does \ml{കിരീടം} mean? (\textquotedblleft{}A crown\textquotedblright{}.) Say it once more.
\end{tcolorbox}
\section[\texorpdfstring{\ml{പട്ടാഭിഷേകം} (paṭṭābhiṣēkaṁ)}{പട്ടാഭിഷേകം (paṭṭābhiṣēkaṁ)}]{\ml{പട്ടാഭിഷേകം} (paṭṭābhiṣēkaṁ) — a coronation}

\label{lesson:ML-C398-pattabhisekam}



\begin{quote}
Before the new one: say the Malayalam for a throne, then the Malayalam for a crown.
\end{quote}

\subsection*{You'll want to know: \ml{പട്ടാഭിഷേകം}}
\textbf{\ml{പട്ടാഭിഷേകം}} — \emph{paṭṭābhiṣēkaṁ} — \textquotedblleft{}a coronation\textquotedblright{}.

\begin{quote}
\ml{പുതിയ} \ml{രാജാവിന്റെ} \ml{പട്ടാഭിഷേകം} \ml{ഒരു} \ml{വലിയ} \ml{ഉത്സവം} \ml{ആയിരുന്നു}.
\end{quote}

\begin{quote}
\emph{putiya rājāvinṟe paṭṭābhiṣēkaṁ oru valiya utsavaṁ āyirunnu.}
\end{quote}

\begin{quote}
The new king's coronation was a great festival.
\end{quote}

\ml{പട്ടം}, a royal band, + \ml{അഭിഷേകം}, an anointing.

\begin{cousinweb}
\textbf{\ml{പട്ടാഭിഷേകം}} (\emph{paṭṭābhiṣēkaṁ}) is built from Sanskrit \emph{paṭṭa}, \textquotedblleft{}cloth; a fillet, a diadem\textquotedblright{} (CDIAL 7700), and \emph{abhiṣeka}, \textquotedblleft{}anointing; the inauguration of a king\textquotedblright{} (Monier-Williams).
\end{cousinweb}

\begin{grammarlens}[title={What you've built}]
One more word for telling what happened.
\end{grammarlens}

\subsection*{Your turn}
\begin{itemize}
\raggedright
  \item \emph{Say it:} \emph{paṭṭābhiṣēkaṁ}
  \item \emph{Say it:} \emph{paṭṭābhiṣēkaṁ}, once more
  \item \emph{Say it:} \emph{putiya rājāvinṟe paṭṭābhiṣēkaṁ oru valiya utsavaṁ āyirunnu.} — The new king's coronation was a great festival.
  \item \emph{Recall:} say the Malayalam for a throne, then the Malayalam for a crown, then say \emph{paṭṭābhiṣēkaṁ} again
\end{itemize}

\begin{tcolorbox}[breakable,colback=teal!4,colframe=teal!35!black,title={Before you move on}]
What does \ml{പട്ടാഭിഷേകം} mean? (\textquotedblleft{}A coronation\textquotedblright{}.) Say it once more.
\end{tcolorbox}
\section[\texorpdfstring{\ml{ഭരണാധികാരി} (bharaṇādhikāri)}{ഭരണാധികാരി (bharaṇādhikāri)}]{\ml{ഭരണാധികാരി} (bharaṇādhikāri) — a sovereign; a ruler}

\label{lesson:ML-C398-bharanadhikari}



\begin{quote}
Before the new one: say the Malayalam for a crown, then the Malayalam for a coronation.
\end{quote}

\subsection*{You'll want to know: \ml{ഭരണാധികാരി}}
\textbf{\ml{ഭരണാധികാരി}} — \emph{bharaṇādhikāri} — \textquotedblleft{}a sovereign; a ruler\textquotedblright{}.

\begin{quote}
\ml{അക്കാലത്ത്} \ml{ആ} \ml{രാജ്യത്തിന്റെ} \ml{ഭരണാധികാരി} \ml{ഒരു} \ml{രാജ്ഞി} \ml{ആയിരുന്നു}.
\end{quote}

\begin{quote}
\emph{akkālattŭ ā rājyattinṟe bharaṇādhikāri oru rājñi āyirunnu.}
\end{quote}

\begin{quote}
In those days the ruler of that country was a queen.
\end{quote}

One who holds the authority to rule a land.

\begin{cousinweb}
\textbf{\ml{ഭരണാധികാരി}} (\emph{bharaṇādhikāri}) is built from Sanskrit \emph{bharaṇa}, \textquotedblleft{}bearing, maintaining\textquotedblright{}, and \emph{adhikārin}, \textquotedblleft{}possessing authority; a governor\textquotedblright{} (Monier-Williams).
\end{cousinweb}

\begin{grammarlens}[title={What you've built}]
One more word for telling what happened.
\end{grammarlens}

\subsection*{Your turn}
\begin{itemize}
\raggedright
  \item \emph{Say it:} \emph{bharaṇādhikāri}
  \item \emph{Say it:} \emph{bharaṇādhikāri}, once more
  \item \emph{Say it:} \emph{akkālattŭ ā rājyattinṟe bharaṇādhikāri oru rājñi āyirunnu.} — In those days the ruler of that country was a queen.
  \item \emph{Recall:} say the Malayalam for a crown, then the Malayalam for a coronation, then say \emph{bharaṇādhikāri} again
\end{itemize}

\begin{tcolorbox}[breakable,colback=teal!4,colframe=teal!35!black,title={Before you move on}]
What does \ml{ഭരണാധികാരി} mean? (\textquotedblleft{}A sovereign; a ruler\textquotedblright{}.) Say it once more.
\end{tcolorbox}
\section[\texorpdfstring{\ml{വാഴുക} (vāḻuka)}{വാഴുക (vāḻuka)}]{\ml{വാഴുക} (vāḻuka) — to reign}

\label{lesson:ML-C398-valuka}



\begin{quote}
Before the new one: say the Malayalam for a coronation, then the Malayalam for a sovereign.
\end{quote}

\subsection*{You'll want to know: \ml{വാഴുക}}
\textbf{\ml{വാഴുക}} — \emph{vāḻuka} — \textquotedblleft{}to reign\textquotedblright{}.

\begin{quote}
\ml{ആ} \ml{രാജാവ്} \ml{നാൽപത്} \ml{വർഷം} \ml{വാണു}.
\end{quote}

\begin{quote}
\emph{ā rājāvŭ nālpatŭ varṣaṁ vāṇu.}
\end{quote}

\begin{quote}
That king reigned for forty years.
\end{quote}

Past: \ml{വാണു}, \emph{vāṇu}.

Literary and old: to reign, to live in glory. The past is \ml{വാണു}.

\begin{cousinweb}
\textbf{\ml{വാഴുക}} (\emph{vāḻuka}) is native Dravidian: DEDR 5372 lists Malayalam \emph{vāḻuka}, \textquotedblleft{}to live, to live well or happily, to rule, to reign\textquotedblright{}.
\end{cousinweb}

\begin{grammarlens}[title={What you've built}]
One more word for telling what happened.
\end{grammarlens}

\subsection*{Your turn}
\begin{itemize}
\raggedright
  \item \emph{Say it:} \emph{vāḻuka}
  \item \emph{Say it:} \emph{vāḻuka}, once more
  \item \emph{Say it:} \emph{ā rājāvŭ nālpatŭ varṣaṁ vāṇu.} — That king reigned for forty years.
  \item \emph{Recall:} say the Malayalam for a coronation, then the Malayalam for a sovereign, then say \emph{vāḻuka} again
\end{itemize}

\begin{tcolorbox}[breakable,colback=teal!4,colframe=teal!35!black,title={Before you move on}]
What does \ml{വാഴുക} mean? (\textquotedblleft{}To reign\textquotedblright{}.) Say it once more.
\end{tcolorbox}
